\section{Stability and Guarantees}
\label{sec:stability}

This section defines the notion of stability that is used, proves that the rule layer has it, and
gives the condition under which the correction cannot reduce contrast. Section~\ref{sec:theory_check}
tests both results on real images.

\subsection{What stability means here}
\label{sec:stability_def}
This is not stability in the sense of Lyapunov and we claim no fixed point. The quantity that must
not behave erratically is the decision of the rule layer, a function of the five rule truth values,
which are exactly what shift when the backbone or the scanner changes.

\begin{definition}[Stability of the allocation]
\label{def:stability}
Let $\mathbf{r} = (r_1, \ldots, r_5) \in [0,1]^5$ be the vector of rule truth values at one pixel and
let $a(\mathbf{r})$ be the allocation of Equation~\eqref{eq:allocation} at that pixel with the
parameters held fixed. The allocation is said to be stable with constant $L$ if
\begin{equation}
\big| a(\mathbf{r}) - a(\mathbf{r}') \big| \;\le\; L \, \lVert \mathbf{r} - \mathbf{r}' \rVert_\infty
\qquad \text{for all } \mathbf{r}, \mathbf{r}' \in [0,1]^5 .
\label{eq:stability_def}
\end{equation}
\end{definition}

In words, a small change in the evidence produces a proportionally small change in the decision, and
$L$ is the worst case ratio between the two. Without such a bound the layer could swing from applying
no correction to applying all of it because one property score moved slightly, which is precisely
the behaviour a clinical user cannot audit.

\subsection{The rule layer is stable, with an exact constant}

\begin{theorem}[Allocation stability]
\label{thm:stability}
The allocation of Equation~\eqref{eq:allocation} satisfies Definition~\ref{def:stability} with
\begin{equation}
L \;=\; C_1 w_1 + C_2 w_2 + C_3 w_3 + C_4 w_4 + C_5 w_5 ,
\label{eq:lipschitz}
\end{equation}
where $w_k = \sigma(\theta_k)$ are the learned rule weights. The constant is attained, so it is the
exact Lipschitz constant and not merely an upper bound. Moreover $L < C_1 + C_2 + C_3 + C_4 + C_5 =
0.75$ for every possible setting of the parameters.
\end{theorem}

The allocation is affine in $\mathbf{r}$ with coefficients $\pm C_k w_k$, so the bound follows from
the triangle inequality and is attained when each $r_k - r'_k$ takes the sign of its coefficient. The
full proof is given in the supplementary file.

\begin{corollary}[Behaviour under a change of backbone]
\label{cor:transfer}
Suppose replacing one backbone by another changes the rule truth values at a pixel by at most
$\epsilon$ in the maximum norm, with the parameters unchanged. Then the allocation at that pixel
changes by at most $L\epsilon$, with $L$ from Equation~\eqref{eq:lipschitz}.
\end{corollary}

This is a statement about the routing rule alone. It says the decision layer degrades gracefully when
the evidence shifts, not that image quality transfers, which is the empirical question of
Section~\ref{sec:results_ood}.

\subsection{When the correction cannot reduce contrast}
Contrast to noise ratio is the measure this method targets most directly. Write
$\mathrm{CNR}(\mathbf{t}) = (\mu_\mathcal{T}(\mathbf{t}) - \mu_\mathcal{B}(\mathbf{t})) / \sigma_\mathcal{B}(\mathbf{t})$
for the separation between tissue and background divided by the variability of the background.

\begin{theorem}[Contrast margin]
\label{thm:cnr}
Let $\eta_\mathcal{T} = \mu_\mathcal{T}(\mathbf{q}) - \mu_\mathcal{T}(\mathbf{b}) \ge 0$ be the gain in
tissue mean and let $\eta_\mathcal{B} \ge \mu_\mathcal{B}(\mathbf{q}) - \mu_\mathcal{B}(\mathbf{b})$
bound the leakage into the background. If the correction does not raise the background variability,
that is $\sigma_\mathcal{B}(\mathbf{q}) \le \sigma_\mathcal{B}(\mathbf{b})$, then the output of
Equation~\eqref{eq:blend} satisfies
\begin{equation}
\mathrm{CNR}(\hat{\mathbf{x}}) \;\ge\;
\frac{\big(\mu_\mathcal{T}(\mathbf{b}) - \mu_\mathcal{B}(\mathbf{b})\big) + w\,(\eta_\mathcal{T} - \eta_\mathcal{B})}
     {\sigma_\mathcal{B}(\mathbf{b})} .
\label{eq:cnr_bound}
\end{equation}
In particular, whenever $\eta_\mathcal{T} > \eta_\mathcal{B}$ the correction cannot reduce the
contrast to noise ratio.
\end{theorem}

\noindent The proof is a convex combination argument for the numerator and Minkowski's
inequality for the denominator. It is given in full in the supplementary file.


Three parts of the design make the hypotheses of Theorem~\ref{thm:cnr} hold in practice. The gate of
Equation~\eqref{eq:gate} keeps $\eta_\mathcal{B}$ small by suppressing change in the dark region, the
brightness branch of the gain network cannot be negative, which keeps $\eta_\mathcal{T}$ non negative,
and the background smoothing lowers $\sigma_\mathcal{B}(\mathbf{q})$ rather than raising it.
Section~\ref{sec:theory_check} measures all three on every test image and reports how often the
hypotheses hold.

Both results are deliberately narrow. Theorem~\ref{thm:stability} bounds how far the routing decision
can move, not whether it is correct, Theorem~\ref{thm:cnr} concerns contrast alone and says nothing
about the other five measures, and neither promises that the output is closer to the clean reference,
which Section~\ref{sec:results_id} shows it usually is not. The guarantees concern the policy, not the
denoiser underneath it.
